% Generated draft: requires mathematical and source review.
\documentclass[11pt]{article}
\usepackage[margin=25mm]{geometry}
\usepackage{fontspec}
\setmainfont{Latin Modern Roman}
\usepackage{amsmath,amssymb,mathtools,unicode-math}
\setmathfont{Latin Modern Math}
\usepackage{enumitem,needspace,xurl,hyperref}
\setlength{\parindent}{0pt}
\setlength{\parskip}{5pt}
\setlength{\emergencystretch}{3em}
\newcommand{\N}{\mathbb{N}}
\newcommand{\Z}{\mathbb{Z}}
\newcommand{\Q}{\mathbb{Q}}
\newcommand{\R}{\mathbb{R}}
\newcommand{\C}{\mathbb{C}}
\newcommand{\F}{\mathbb{F}}
\newcommand{\A}{\mathbb{A}}
\newcommand{\PP}{\mathbb{P}}
\newcommand{\E}{\mathbb{E}}
\newcommand{\eps}{\varepsilon}
\newcommand{\dd}{\,\mathrm{d}}
\begin{document}
\section*{UnsolvedMath ID: 800720}
\textbf{Source problem code:} catalogue entry
\subsection{Definitions and mathematical statement}
Let \(F\) be a field. A central division \(F\)-algebra of degree \(3\) is a division algebra \(D\) with center \(F\) and \(\dim_F D=9\). A field extension \(L/F\) splits \(D\) if
\[
D\otimes_F L\cong M_3(L)
\]
as \(L\)-algebras. Write \(\operatorname{Br}(F)\) for the Brauer group of finite-dimensional central simple \(F\)-algebras, \([D]\) for the class of \(D\), and
\[
\langle[D]\rangle=\{n[D]:n\in\Z\}
\]
for its generated subgroup.

Under the standard interpretation that the two division algebras are central over a common field, and the literal interpretation that the splitting fields range over all field extensions of the specified transcendence degree, the question is: for every field \(F\) and every pair of central division \(F\)-algebras \(D_1,D_2\) of degree \(3\), does
\[
\begin{gathered}
\text{for every field extension }L/F\text{ with }\operatorname{trdeg}_F L\leq 1,\\
D_1\otimes_F L\cong M_3(L)
\quad\Longleftrightarrow\quad
D_2\otimes_F L\cong M_3(L)
\end{gathered}
\]
imply
\[
\langle[D_1]\rangle=\langle[D_2]\rangle
\quad\text{in }\operatorname{Br}(F)?
\]
The transcendence-degree bound includes algebraic extensions. The supplied excerpt does not explicitly specify centrality, the common base field, characteristic restrictions, or whether only finitely generated extensions are intended. The displayed formulation makes its interpretation explicit; these points require comparison with the full source context. The source describes this as what ``appears to be'' the smallest open case; that qualification and the registry's open label are supplied assessments, not independently verified conclusions.
\subsection{Short English statement}
Must two central division algebras of degree \(3\) over a common field generate the same subgroup of its Brauer group if exactly the same extensions of transcendence degree at most \(1\) split them? The interpretation of the base field and the permitted extensions requires source review.
\subsection{Sources}
\begin{itemize}
\item Local UnsolvedMath registry, permanent ID \(800720\), exact source problem code \texttt{catalogue entry}, title \emph{Open Problems on Central Simple Algebras --- Problem 34}. The record attributes the problem to Auel, Brussel, Garibaldi, and Vishne, with proposed year \(2010\), and labels it open.
\item Auel, Brussel, Garibaldi, and Vishne, \emph{Open Problems on Central Simple Algebras} (2010), Problem 34. Supplied source URL: \url{https://arxiv.org/abs/1006.3304}. The supplied background records access on 2026-10-16 and extraction by \texttt{arxiv\_latex}; these are source-provenance statements and do not indicate independent access or verification for this draft.
\item The supplied excerpt qualifies the degree-\(3\) question as what ``appears to be'' the smallest open case. Its open status and this comparison are retained as supplied source assessments, without a new literature assessment.
\end{itemize}
\end{document}
